\chapter*{Anot\=acija (LV)}
\addcontentsline{toc}{chapter}{Anot\=acija (LV)}

\textbf{Nosaukums:} M\=ak\-sl\=ig\=a intelekta Telegram bota izstr\=ade person\-aliz\=etiem d\=avanu ieteikumiem, izman\-tojot prompt inženierijas tehnikas.

D\=avanas izv\=ele ir ikdienišks, bet laikietilp\=igs uz\-de\-vums. Tieš\-saistes veikali r\=ada daudz pre\-\v cu, bet ne\-pied\=av\=a person\-aliz\=etus padomus. Mo\-derni lielie va\-lo\-das mode\-\c li (LLM) var ana\-liz\=et d\=avanas sa\-\c n\=em\=eja profilu un ie\-te\-ikt da\-v\=anu ide\-jas da\-bisk\=a va\-lo\-d\=a. Ta\-\v cu LLM at\-bilžu kva\-lit\=ate ir \c lo\-ti at\-ka\-r\=iga no ievades teksta (prompta) di\-zaina. Šo problē\-mu p\=et\=i jauna pē\-t\=ij\=umu jo\-ma — prompt in\-žen\-ierija.

Darba m\=er\c kis ir iz\-strā\-d\=at Tele\-gram botu, kas iz\-manto LLM (Grok no xAI), lai sniegtu person\-aliz\=etus d\=avanu iete\-ikumus, un sa\-l\=id\-zin\=at \v cetras prompt in\v zen\-ierijas stra\-t\=e\-\v gi\-jas šim uz\-de\-vu\-mam: Naive, Cons\-trai\-ned, Chain-of-Thought (CoT) un Per\-sona-based.

Bots tika izstr\=ad\=ats Python valod\=a, izmantojot aiogram bibliot\=eku un SQLite datu b\=azi. Lieto\-t\=ajs iziet \=isu anketu ar sešiem strukturi\-z\=etiem jaut\=ajumiem un sa\-\c nem trīs personaliz\=etas d\=avanu idejas. Sis\-t\=ema ir izvietota PythonAnywhere m\=ako\-\c na pla\-tform\=a.

Tika veikts kontrol\=ets eksperiments. Tika sagatavotas 30 sin\-t\=etiskas lietot\=aju per\-sonas. Visas \v cetras stra\-t\=e\-\v gi\-jas tika izpild\=itas uz vis\=am person\=am, kop\=a — 120 testa gad\=ijumi. Katra atbilde tika nov\=ert\=eta p\=ec piec\=am metrik\=am: re\-le\-vance, ra\-doš\=ums, specifis\-kums, halluci\-n\=acijas iezīme un at\-bildes laiks. Bots tika test\=ets ar\=i ar 8 re\=alajiem lietot\=ajiem.

Rezult\=ati apstiprina hipot\=ezi, ka uzlabot\=as stra\-t\=e\-\v gi\-jas dod ie\-v\=ero\-jami labākus iete\-ikumus nek\=a vienkārš\=a (naive) pieeja. Lab\=ak\=a kop\=ej\=a stra\-t\=e\-\v gi\-ja bija Chain-of-Thought (relevance 4.43/5, hallucin\=acija 8.9\%). Visradoš\=ak\=a stra\-t\=e\-\v gi\-ja bija Persona-based (radoš\=ums 4.53/5). Visas \v cetras stra\-t\=e\-\v gi\-jas atbild\=eja mazāk nek\=a 10 sekund\=es. 75\% lietot\=aju teica, ka izmantotu botu atk\=al.

Darb\=a ir pied\=av\=ati turp\-m\=akie uzla\-bo\-jumi: integr\=acija ar p\=ardošanas platform\=am, daudz\-mod\=ala ievade, daudzva\-lod\=u atbalsts un hib\-r\=id\=a prompt pipeline.

Darbs ap\-stiprina, ka prompt in\-žen\-ierija ir re\=ala un izm\=erojama in\-žen\-ierijas dar\-b\=iba. Pied\=av\=at\=a meto\-do\-lo\-\v gi\-ja ir do\-m\=ena ne\-at\-ka\-r\=iga un to var izmantot citos personaliz\=etos rekomend\=aciju uzde\-vumos.

\textbf{Apj\=oms:} 55 lappuses, 14 tabulas, 7 at\-tē\-li, zx atsauces, 22 pielikumi.
